\scanpage{18}\mspage{3}
\[
\mathrm{p}_{*}\,\Omega^{\mathrm{reg}}_{\mathrm{E}/\mathrm{S}}\handwritten{\xrightarrow{\sim}}\mathrm{o}^{*}\,\Omega^{1}_{\mathrm{E}/\mathrm{S}}=\omega\quad,
\]
et que pour la flèche adjointe à la flèche inverse
\[
\mathrm{p}^{*}\,\omega\handwritten{\xrightarrow{\sim}}\Omega^{\mathrm{reg}}_{\mathrm{E}/\mathrm{S}}\qquad.
\]
\margin{\handwritten{?}}
Les sections de $\Omega^{\mathrm{reg}}_{\mathrm{E}/\mathrm{S}}$ ainsi associées aux sections de $\omega$ s'appellent les différentielles invariantes.

Soit $\unsure{\grave{\pi}|\dot{\pi}|\pi}$ une section inversible de $\omega$ ; il existe, localement pour la topologie de Zariski, une base $(\mathrm{y},\mathrm{x},1)$ de $\mathrm{p}_{*}\handwritten{\mathcal{O}}(3\mathrm{e})$ telle que l'image de $\mathrm{y}$ dans $\omega^{\handwritten{\otimes}-3}$ soit $-\pi^{\handwritten{\otimes}-3}$, que $\mathrm{x}$ appartienne à $\mathrm{p}_{*}\handwritten{\mathcal{O}}(2\mathrm{e})$, et que son image dans $\omega^{\handwritten{\otimes}-2}$ soit $\pi^{\handwritten{\unsure{\otimes|\circ}}-2}$. Si $\pi'$ est une autre section inversible de $\omega$, les bases correspondantes s'écrivent toutes, sous forme implicite (pour $\mathrm{r},\mathrm{s},\mathrm{t}$ convenables)
\[
\handwritten{\left\{}\begin{aligned}
\mathrm{x}&=\mathrm{u}^{2}\,\mathrm{x}'+\mathrm{r}\\
\mathrm{y}&=\mathrm{u}^{3}\,\mathrm{y}'+\mathrm{s}\,\mathrm{u}^{2}\,\mathrm{x}'+\mathrm{t}\\
\pi'&=\mathrm{u}\,\pi\qquad.
\end{aligned}\right.
\tag*{(1.2)}
\]

Le faisceau\handwritten{$\backslash$} $\handwritten{\mathcal{O}}(3\mathrm{e})$ est très ample et $(\mathrm{x},\mathrm{y},\tikz[baseline=(digit.base)]{\node[inner sep=0pt] (digit) {$1)$};\handwritten{\draw[thin,gray,rotate around={-16:(digit.center)}] (digit.center) ellipse [x radius=.68em,y radius=1.35em];}}$ sont les coordonnées projectives d'un plongement de $\mathrm{E}$ dans $\mathrm{P}_{2}/\mathrm{S}$.

La courbe image vérifiera une équation
\[
\mathrm{y}^{2}+\mathrm{a}_{1}\,\mathrm{x}\,\mathrm{y}+\mathrm{a}_{3}\,\mathrm{y}
=\mathrm{x}^{3}+\mathrm{a}_{2}\,\mathrm{x}^{2}+\mathrm{a}_{4}\,\mathrm{x}+\mathrm{a}_{6}\qquad.
\tag*{(1.3)}
\]

Pour vérifier l'existence et l'unicité de cette équation, on
\clearpage
% archive sheet number: 17
